% ============================================================
%  Setmana 33 — Unitat 6 · Angles i distàncies a l'espai
%  Exercicis: 90, 92, 97, 115, 139
% ============================================================

\setmana{33}{Unitat 6 \textperiodcentered{} Angles i distàncies a l'espai}{2 de maig de 2027}%
        {90, 92, 97, 115 i 139}

\begin{errada}{Exercici 90 a)}
  \ondiu{$r:\ \dfrac{x}{1}=\dfrac{y+1}{-1}=\dfrac{z+1}{-1}$}
  \hadedir{$r:\ \dfrac{x}{1}=\dfrac{y+1}{-1}=\dfrac{z-1}{-1}$}
  \motiu{La recta passa per $A(0,-1,1)$, on $z=1$. La línia següent,
    $x=-z+1$, ja és correcta.}
\end{errada}

\begin{errada}{Exercici 90 b)}
  \ondiu{$d(A,\pi)=\dfrac{|3\cdot(-2)+5\cdot 5-2\cdot(-4)+7|}
    {\sqrt{3^2+5^2+(-2)^2}}=\dfrac{34}{\sqrt{38}}$ u}
  \hadedir{$d(A,\pi)=\dfrac{|3\cdot(-2)+5\cdot 1-2\cdot(-4)+7|}
    {\sqrt{3^2+5^2+(-2)^2}}=\dfrac{14}{\sqrt{38}}$ u}
  \motiu{La segona coordenada del punt $A(-2,1,-4)$ és $1$, no $5$. La
    distància és $\tfrac{14}{\sqrt{38}}=\tfrac{7\sqrt{38}}{19}\approx 2{,}27$ u,
    que és també la distància de $A$ al punt $Q$ que s'ha trobat.}
\end{errada}

\begin{errada}{Exercici 139 a)}
  \ondiu{$r:\ \dfrac{x-1}{1}=\dfrac{y-2}{2}=\dfrac{z-1}{3}$}
  \hadedir{$r:\ \dfrac{x-1}{1}=\dfrac{y-2}{1}=\dfrac{z-1}{3}$}
  \motiu{El vector director de la recta és el normal del pla,
    $\vec{n}=(1,1,3)$: el segon denominador és $1$.}
\end{errada}

\begin{errada}{Exercici 139 b)}
  \ondiu{$Q\left(\tfrac{1}{2},\,1,\,-\tfrac{1}{2}\right)$}
  \hadedir{$Q\left(\tfrac{5}{11},\,\tfrac{16}{11},\,-\tfrac{7}{11}\right)$}
  \motiu{Amb la recta correcta, $(1+t,\;2+t,\;1+3t)$, el pla dona $6+11t=0$,
    és a dir, $t=-\tfrac{6}{11}$.}
\end{errada}

\begin{errada}{Exercici 139 c)}
  \ondiu{$d(P,\pi)=\sqrt{\dfrac{7}{2}}$ u}
  \hadedir{$d(P,\pi)=\dfrac{6}{\sqrt{11}}=\dfrac{6\sqrt{11}}{11}$ u}
  \motiu{Amb la fórmula, $d=\dfrac{|1+2+3\cdot 1|}{\sqrt{1+1+9}}=
    \dfrac{6}{\sqrt{11}}\approx 1{,}81$ u. El valor de la solució és la
    distància de $P$ a un punt del pla que no és el peu de la perpendicular.}
\end{errada}

%\senseerrades{92, 97 i 115}
